% Chapter 13: John Pardon — Geometry, Topology, and Knots

\chapter{John Pardon: Knot Distortion}

\section{The Problem}

\subsection{Informal}

For a rectifiable embedding $\gamma:S^1\to\mathbb{R}^3$, compare intrinsic arclength along the curve with Euclidean distance in space. Its distortion is
\[
\delta(\gamma)=\sup_{p\ne q\in\gamma}\frac{d_\gamma(p,q)}{\lVert p-q\rVert}.
\]
For a knot type $K$, define $\delta(K)$ as the infimum of $\delta(\gamma)$ over all representatives of $K$.

\subsection{Formal}

Gromov asked whether every knot type has a representative with distortion bounded by a universal constant, and in particular whether the distortion of torus knots remains bounded as their parameters grow.

\subsection{Historical Context}

John Pardon answered the torus-knot form of the question in work published in 2011. His argument gives a quantitative lower bound for knots lying on embedded surfaces; it does not use an asymptotic Conley--Zehnder capacity.

\subsection{The Basic Example}

For a round circle, the maximum ratio is $\pi/2$, since intrinsic distance is half the circumference while the corresponding chord has length twice the radius. Thus the unknot has distortion at most $\pi/2$ (and this is the value for a round representative), not 1.

\section{Mathematical Theory}

\subsection{Curves on Embedded Surfaces}

Let $F\subset\mathbb{R}^3$ be a PL embedded closed surface of genus at least one. If $\beta$ is an essential simple loop on $F$, let $K_\beta$ denote the corresponding knot and let $I(F,\beta)$ be its minimum geometric intersection number with a loop on $F$ bounding a disk whose interior misses $F$.

The key point is that intersection number is topological, while distortion is metric. Pardon proves that sufficiently complicated intersection data forces every representative of the knot type to have large distortion.

\section{The Solution}

\begin{theorem}[Pardon, 2011 \cite{pardon2011}]
For a knot $K_\beta$ represented by an essential loop $\beta$ on an embedded surface $F$,
\[
\delta(K_\beta)\geq \frac{1}{160}I(F,\beta).
\]
In particular, for the $(p,q)$ torus knot,
\[
\delta(T_{p,q})\geq \frac{1}{160}\min(p,q).
\]
\end{theorem}

Taking $p=q\to\infty$ shows that knot distortion is unbounded across knot types. The proof converts intersections on the surface into pairs of points that must be far apart in the knot metric but close in the ambient geometry, and then controls how this obstruction survives under arbitrary re-embedding.

\subsection{What the Theorem Does Not Say}

The theorem gives lower bounds for explicit families of knot types. It does not determine the optimal distortion of every individual knot, nor does it introduce a universal contact invariant. Questions about sharp constants and higher-dimensional analogues remain active.

\section{Impact}

Pardon's result answered a question posed by Gromov and supplied a robust geometric method for proving lower bounds on knot distortion. It also illustrates a useful general strategy: topological intersection data can force metric complexity.

\section{Exercises}

\begin{exercise}[Basic]
Compute the distortion of a round circle using intrinsic arclength and Euclidean chord distance.
\end{exercise}

\begin{exercise}[Intermediate]
Explain why a lower bound growing with $\min(p,q)$ proves that distortion is unbounded over the family of torus knots.
\end{exercise}

\section{Timeline}

\begin{tabularx}{\linewidth}{lX}
\toprule
\textbf{Year} & \textbf{Event} \\ \midrule
1983 & Gromov poses the knot-distortion question \\
2010 & Pardon posts the preprint on distortion of knots on embedded surfaces \\
2011 & The result appears in the \emph{Annals of Mathematics} \\
\bottomrule
\end{tabularx}

\section{Citations}

The primary reference is Pardon’s paper listed as \cite{pardon2011} in the bibliography.
